% ЗӨВХӨН БҮТЭН НООРОГТ (main.tex) хэвлэгдэнэ. Багшид үзүүлэх ноорогт (main_teacher.tex) орохгүй.
% 2.1-ийн төгсгөл: их сургуулийн жишээ (Marcel; 2026.11.16 (4)-нд хэрэглэгч багшийн нооргоос авсан) ба дэд хэсэг «Хариулт байхгүй асуулт».
Ийм аргыг их сургуулийн туслах системд хэрэглэсэн жишээ бий. Германы нэгэн их сургуулийн Marcel системд ажилтнуудтай хамт бэлтгэж, холбогдох баримттай нь холбосон 36 бэлэн асуулт–хариултаас хайдаг хэсэг зөв баримтыг ерөнхий хайлтаас дээгүүр эрэмбэлсэн \cite{trienes2025marcel}. Зөв баримт жагсаалтын хэддүгээр байрт гарсны урвуу утгын дундаж (\acs{MRR}) энэ хэсэгт 0.70, ерөнхий хайлтад 0.40 байв. Харин Marcel-д хариултыг хэлний загвар бичдэг тул энэ дүн эх сурвалжийн өгүүлбэрийг үгчлэн харуулах аргыг шууд батлахгүй.

\subsection{Хариулт байхгүй асуулт}
\label{sec:lit-unanswerable}
% SQuAD 2.0 (RAG08): агент бүтэн уншиж papers/16, papers/22-т бүртгэсэн. Хэрэглэгч W03-т уншина.

Сургуулийн туслахад хүн мэдээллийн санд байхгүй зүйлийг ч асууна. Ийм үед систем таамаглаж хариулах нь аюултай. Rajpurkar нар хариулт нь өгөгдсөн бичвэрт байхгүй 53 мянга гаруй асуулттай SQuAD 2.0 өгөгдлийг гаргасан \cite{rajpurkar2018squad2}. Асуултууд тухайн бичвэртэй холбоотой, үнэмшилтэй харагддаг боловч хариулт нь бичвэрт байхгүй. Өмнөх SQuAD 1.1 өгөгдөл дээр F1 (нарийвчлал ба бүрэн байдлыг нэгтгэсэн оноо) 85.8 авдаг загварын бүтэц хариулт байхгүй асуулт нэмэгдсэн SQuAD 2.0 дээр 66.3 авсан бол хүн 89.5 авсан \cite{rajpurkar2018squad2}. Хэрэглэгчийн зорилгыг ангилдаг системүүд ч хамрах хүрээндээ байгаа асуултыг сайн ангилдаг боловч хүрээнээс гадуурх асуултыг танихдаа сул байдаг \cite{larson2019clinc}. Marcel системд зөв баримтыг өгсөн үед хариулах боломжгүй асуултаас зөв татгалзах нь 0.68-аас 0.93 болж сайжирсан \cite{trienes2025marcel}.

Эдгээр ажил хариулахаас татгалзахыг хэмжсэн бөгөөд даалгавар, хэл нь энэ ажлынхаас өөр. Энэ ажилд систем татгалзаж зогсохгүй, хандах алба, өрөөг заах шаардлагатай. Иймд цуглуулах асуултаас хариулт нь сургуулийн мэдээлэлд байхгүй асуултыг тусад нь тэмдэглэж, ийм үед систем хандах алба, өрөөг зөв зааж байгаа эсэхийг шалгана (\ref{sec:exp-tree}-р хэсэг).
